\chapter{Databases and the head of a ranked list}
\label{ch:databases}

% Author's note (verification and reproduction): session screenlit-q3
% (287 MB; loaded once, for the matrix and the table) and screenlit-q3-top;
% queries q3a-q3h, q3x, q3y; check by meaning checks/q3*.json; samples
% tier2-sample-q3.tsv (seed 20261126) and sample60-q3.txt (seed
% 20261123); see the block at the top of search.tex.  Sources:
% doc/bibliography.bib, keys Bramer2017 Bramer2016 Aagaard2016 Gehanno2013
% Gusenbauer2020 Haddaway2015 Sampson2006 Heck2024 Wanyama2021 Rajit2025
% Ali2022 Hartling2016.

The procedure of \cref{ch:procedure} searches every database Scholar
offers that answers, and screens the first 100 places of each answer ordered by
relevance.
Both choices assume something about databases: that one or two would not
do, and that the head of a ranked list is the right place to start but
not a place to stop.
This chapter asks the question in the form that would make both choices
unnecessary: \emph{do a few bibliographic databases, or the first
results of a ranked source, find nearly all the relevant records a
review needs?}
\enquote{Nearly all} is taken as a recall that a review could report
without a caveat.

\section{Method}
\label{sec:databases-method}

The search is the one described in \cref{ch:lit-search}.
The third question's queries are listed below, and
\cref{tab:databases-matrix} gives their hits per database.
The support queries are q3a to q3f, the counter queries q3x and q3y; q3g
was added after the check by meaning, with the words of studies that
compare what different databases find, and q3h because two of the named
studies had not come back on the question's concepts.
\begin{description}[style=nextline,font=\normalfont\bfseries]
\item[q3a, support query] {\footnotesize\ttfamily\raggedright (\textquotedbl{}Google Scholar\textquotedbl{} OR Scopus OR \textquotedbl{}Web of Science\textquotedbl{}) AND (coverage OR overlap OR \textquotedbl{}unique records\textquotedbl{} OR \textquotedbl{}unique references\textquotedbl{}) AND (\textquotedbl{}systematic review\textquotedbl{} OR \textquotedbl{}bibliographic databases\textquotedbl{})\par}
\item[q3b, support query] {\footnotesize\ttfamily\raggedright (OpenAlex OR \textquotedbl{}Semantic Scholar\textquotedbl{}) AND (coverage OR overlap OR comparison OR completeness) AND (database OR bibliographic)\par}
\item[q3c, support query] {\footnotesize\ttfamily\raggedright (\textquotedbl{}Google Scholar\textquotedbl{}) AND (\textquotedbl{}systematic review\textquotedbl{} OR \textquotedbl{}evidence review\textquotedbl{}) AND (\textquotedbl{}first 100\textquotedbl{} OR \textquotedbl{}first 200\textquotedbl{} OR \textquotedbl{}first 50\textquotedbl{} OR \textquotedbl{}top 100\textquotedbl{} OR \textquotedbl{}first pages\textquotedbl{} OR \textquotedbl{}first 1000\textquotedbl{})\par}
\item[q3d, support query] {\footnotesize\ttfamily\raggedright (\textquotedbl{}optimal combination\textquotedbl{} OR \textquotedbl{}database combination\textquotedbl{} OR \textquotedbl{}number of databases\textquotedbl{} OR \textquotedbl{}which databases\textquotedbl{} OR \textquotedbl{}additional database\textquotedbl{}) AND (\textquotedbl{}systematic review\textquotedbl{} OR \textquotedbl{}literature search\textquotedbl{}) AND (recall OR unique OR coverage)\par}
\item[q3e, support query] {\footnotesize\ttfamily\raggedright (\textquotedbl{}software engineering\textquotedbl{} OR \textquotedbl{}computer science\textquotedbl{} OR computing OR education OR \textquotedbl{}educational research\textquotedbl{}) AND (\textquotedbl{}digital libraries\textquotedbl{} OR \textquotedbl{}digital library\textquotedbl{} OR \textquotedbl{}bibliographic databases\textquotedbl{} OR \textquotedbl{}search engines\textquotedbl{} OR Scopus OR \textquotedbl{}Web of Science\textquotedbl{} OR \textquotedbl{}Google Scholar\textquotedbl{}) AND (\textquotedbl{}systematic literature review\textquotedbl{} OR \textquotedbl{}systematic review\textquotedbl{} OR \textquotedbl{}literature search\textquotedbl{}) AND (coverage OR overlap OR recall OR \textquotedbl{}search strategy\textquotedbl{} OR completeness)\par}
\item[q3f, support query] {\footnotesize\ttfamily\raggedright (\textquotedbl{}relevance ranking\textquotedbl{} OR \textquotedbl{}relevance-ranked\textquotedbl{} OR \textquotedbl{}first 100\textquotedbl{} OR \textquotedbl{}first 200\textquotedbl{} OR \textquotedbl{}first 300\textquotedbl{} OR \textquotedbl{}top results\textquotedbl{} OR \textquotedbl{}first pages\textquotedbl{}) AND (\textquotedbl{}Google Scholar\textquotedbl{} OR \textquotedbl{}search engine\textquotedbl{} OR \textquotedbl{}web search\textquotedbl{} OR \textquotedbl{}academic search engine\textquotedbl{}) AND (\textquotedbl{}systematic review\textquotedbl{} OR \textquotedbl{}literature search\textquotedbl{} OR \textquotedbl{}grey literature\textquotedbl{})\par}
\item[q3g, support query] {\footnotesize\ttfamily\raggedright (Scopus OR \textquotedbl{}Web of Science\textquotedbl{} OR \textquotedbl{}Google Scholar\textquotedbl{} OR EBSCO OR \textquotedbl{}bibliographic databases\textquotedbl{}) AND (\textquotedbl{}systematic literature review\textquotedbl{} OR \textquotedbl{}systematic review\textquotedbl{} OR \textquotedbl{}systematic reviews\textquotedbl{} OR \textquotedbl{}meta-analyses\textquotedbl{}) AND (\textquotedbl{}common to\textquotedbl{} OR shared OR variation OR \textquotedbl{}choice of\textquotedbl{} OR \textquotedbl{}stand-alone\textquotedbl{} OR \textquotedbl{}single database\textquotedbl{} OR sufficient)\par}
\item[q3h, support query] {\footnotesize\ttfamily\raggedright (\textquotedbl{}multiple databases\textquotedbl{} OR \textquotedbl{}number of databases\textquotedbl{} OR \textquotedbl{}additional databases\textquotedbl{} OR \textquotedbl{}selective database searching\textquotedbl{} OR \textquotedbl{}diminishing returns\textquotedbl{} OR \textquotedbl{}added value\textquotedbl{}) AND (\textquotedbl{}systematic review\textquotedbl{} OR \textquotedbl{}systematic reviews\textquotedbl{} OR \textquotedbl{}meta-analysis\textquotedbl{}) AND (searching OR search OR retrieval OR recall)\par}
\item[q3x, counter query] {\footnotesize\ttfamily\raggedright (\textquotedbl{}Google Scholar\textquotedbl{}) AND (inadequate OR \textquotedbl{}not suitable\textquotedbl{} OR limitations OR \textquotedbl{}not recommended\textquotedbl{} OR \textquotedbl{}cannot be used\textquotedbl{}) AND (\textquotedbl{}systematic review\textquotedbl{} OR searching)\par}
\item[q3y, counter query] {\footnotesize\ttfamily\raggedright (\textquotedbl{}Google Scholar\textquotedbl{} OR Scopus OR \textquotedbl{}Web of Science\textquotedbl{} OR OpenAlex OR \textquotedbl{}Semantic Scholar\textquotedbl{} OR \textquotedbl{}bibliographic databases\textquotedbl{}) AND (coverage OR overlap OR \textquotedbl{}unique records\textquotedbl{} OR recall) AND (missed OR missing OR incomplete OR insufficient OR \textquotedbl{}not sufficient\textquotedbl{} OR \textquotedbl{}not indexed\textquotedbl{}) AND (\textquotedbl{}systematic review\textquotedbl{} OR \textquotedbl{}systematic reviews\textquotedbl{} OR \textquotedbl{}literature search\textquotedbl{})\par}
\end{description}


\begin{table}
  \begin{sidecaption}{%
      Hits per query and database for the third question (databases and
      ranked lists), searched 2026-11-23 with Scholar, as \texttt{scholar
      sessions show} reports them; the databases and their syntax as in
      \cref{tab:stopping-matrix}.
      \enquote{of}: the database reported more records than it returned
      (\enquote{of about}: Semantic Scholar's approximate total).
      \textsuperscript{a}~Incomplete; Scholar recorded \enquote{rate
      limited or HTTP/parse error}, and the search's log showed rate
      limits (HTTP 429).
      \textsuperscript{b}~Incomplete; Scopus stopped with HTTP 400 after
      4,991 records.
      The OpenAlex answer to q3g holds 9,978 of the 26,809 records
      OpenAlex reported, and Scholar did not mark it incomplete.
    }[tab:databases-matrix]
    \footnotesize
    \begin{tabular}{@{}l*{6}{>{\raggedleft\arraybackslash}p{1.45cm}}@{}}
\toprule
Query & OA & S2 & Scopus & WoS & IEEE & arXiv \\
\midrule
q3a & 6,237 & 999 of about 5,207\textsuperscript{a} & 4,266 of 4,272 & 3,574 of 3,576 & 50 & 11 \\
q3b & 581 & 227 of about 228 & 118 & 103 & 6 & 0 \\
q3c & 204 & 104 & 101 & 57 & 0 & 0 \\
q3d & 43 & 59 & 43 & 19 & 0 & 0 \\
q3e & 3,973 & 2,971 of about 3,008 & 1,721 of 1,723 & 1,084 of 1,086 & 149 & 51 \\
q3f & 477 of 479 & 116 & 97 & 59 & 0 & 0 \\
q3g & 9,978 of 26,809 & 7,996 of about 47,668\textsuperscript{a} & 4,991 of 18,141\textsuperscript{b} & 14,448 of 14,458 & 172 & 34 \\
q3h & 6,825 of 6,838 & 999 of about 5,845\textsuperscript{a} & 3,867 of 3,870 & 4,519 of 4,521 & 21 & 2 \\
q3x & 8,296 & 9,956 of about 10,680\textsuperscript{a} & 4,924 of 4,927 & 6,178 of 6,181 & 107 & 71 \\
q3y & 1,001 & 1,279 of about 1,294 & 391 & 374 of 375 & 9 & 4 \\
\bottomrule
\end{tabular}

  \end{sidecaption}
\end{table}

Five studies were named before the search: two on database combinations,
one on the contribution of databases to the results of reviews, one on
Google Scholar in evidence reviews, and one older study of searching
several databases.
The study of the contribution of databases came back only through the
second question's queries, and the older study through no query; it is
not cited.
The check by meaning returned 29 works no session held; the repair query
returned one of them, and 26 comparisons of bibliographic coverage stayed
outside (\cref{sec:search-check}).
Of the 61,348 distinct records, 1,003 were screened in the first tier and
17 more that five or more databases had returned; 58 bore on the question
(22 supporting the claim, 36 qualifying or contradicting it).
Of the 60,268 records left unscreened, none among a random sample of 60
bore on the question (95\,\% interval 0 to 6.0\,\%), and none of a second
sample of 60 from the records that three or more databases had returned.
The records judged bearing are listed in \cref{sec:databases-hits}.

\section{Results}
\label{sec:databases-results}

\paragraph{No single database's search, and few added ones, suffice in health.}
In 58 health reviews, \textcquote[abstract]{Bramer2017}{\textins*{s}ixteen
percent of the included references (291 articles) were only found in a
single database; Embase produced the most unique references (n =
132)}.
Embase, MEDLINE, Web of Science and the first 200 results of Google
Scholar together reached a recall of 98.3\,\%, and the authors estimate
that 60\,\% of published reviews miss more than 5\,\% of the relevant
references because of the databases they searched.
In 23 reviews of musculoskeletal disorders, more databases added little:
\textcquote[abstract]{Aagaard2016}{\textins*{c}umulative median recall for
combined searching in MEDLINE, EMBASE and CENTRAL was 88.9\,\% and
increased to 90.9\,\% when adding 10 additional databases}.

\paragraph{Holding a record is not returning it.}
All 738 studies included in 29 medical reviews were in Google Scholar
when looked up by title: \textcquote[abstract]{Gehanno2013}{All the 738
original studies included in the gold standard database were retrieved
in GS (100\,\%).}
In 120 reviews, Google Scholar held 97.2\,\% of the included references,
but a search returned far fewer near the top:
\textcquote[abstract]{Bramer2016}{\textins*{o}nly 46.4\,\% of the included
references were among the downloadable first 1000 references in GS}.
A test of 28 search systems, read only in its abstract, concludes that
\textcquote[abstract]{Gusenbauer2020}{our findings demonstrate why Google
Scholar is inappropriate as principal search system}.
For OpenAlex and Semantic Scholar the one study found measured holding,
not returning: it looked the articles a health guideline included up by
identifier, then by title.
Read only in its abstract, it reports that \textcquote[abstract]{Rajit2025}{OpenAlex
had the best coverage (98.6\,\%), followed by Semantic Scholar (98.3\,\%),
Embase (96.8\,\%), and PubMed (93.0\,\%)}.

\paragraph{The head of a ranked list.}
A search engine ranked the Boolean results of nine Cochrane reviews by
relevance; the recall of included studies within the first 500 places
was 0.70, and \textcquote[abstract]{Sampson2006}{\textins*{t}he relevance
ranking provided by the search engine was better than expected by
chance}, though it could not tell apart records that reviewers had
already pre-screened.
In environmental science, Google Scholar missed important literature in
five of six case studies, and the authors recommend reading further than
the first page: \textcquote[abstract]{Haddaway2015}{If used in systematic
reviews for grey literature, we recommend that searches of article titles
focus on the first 200 to 300 results.}
The study of Google Search in \cref{ch:stopping} found relevant studies at
places 215 and 227.

\paragraph{Outside health.}
\begingroup\sloppy
In three reviews in education, \textcquote[abstract]{Heck2024}{discipline-\hspace{0pt}specific
databases outperform international multidisciplinary sources, and a
combination of discipline-\hspace{0pt}specific international and national sources is
most efficient in finding a high proportion of relevant literature}.
In a social-science review, read only in its abstract,
\textcquote[abstract]{Wanyama2021}{only 130 articles (5.7\,\% of 2267
retrieved) were found common to all three databases}.
For reference and citation searching in software engineering, one case
(\textcquote[abstract]{Ali2022}{\textins*{i}n this particular case}) read
only in its abstract suggests that \textcquote[abstract]{Ali2022}{Scopus
seems to be more effective and useful for snowball sampling than GS for systematic
secondary studies attempting to identify peer-reviewed literature}.
How much the database choice changes a review's results, rather than its
records, is in \cref{ch:abbreviated}; the cross-sectional study of 129
Cochrane reviews there \parencite{Hartling2016} is also one of this
question's named studies.
\par\endgroup

\section{Discussion and limitations}
\label{sec:databases-discussion}

The recall figures are from health research and from databases Scholar
does not search (Embase, MEDLINE, CENTRAL, Google Scholar); they show that
no single source's search sufficed there, not which combination suffices in
computing or education.
In education the discipline's own databases did best, and Scholar
searches none of them; a claim about education should say so.
No study found tests the keyword or ranked retrieval of OpenAlex or
Semantic Scholar; the one that measures them counts records held by
identifier, and was read only in its abstract.
Our own measurements are the only evidence here about the head of
OpenAlex, arXiv and IEEE Xplore lists, and they are one chapter's
sessions (\cref{sec:measured-chapter}).
The search's own answers were large and partly incomplete
(\cref{tab:databases-matrix}), so for this question the first tier and the
sample stand for more unscreened records than for the others.

\section{Conclusion}
\label{sec:databases-conclusion}

No: in health research no single database's search found what the
reviews included, three or four together came close, and further ones
added
little; holding a record is not returning it, since a ranked source that
held almost every included study returned fewer than half of them in its
first 1,000 places; and outside health the right databases differ.
For the procedure this means searching every database that answers,
reading the head of each ranked list first without stopping there, and
not carrying the health figures over to computing or education.

\section{Future work}
\label{sec:databases-future}

This search did not obtain the full texts of the test of 28 search
systems, of the study of OpenAlex and Semantic Scholar, of the
social-science and software-engineering comparisons, or the older study
of searching several databases.
The field lacks a measurement of the recall of OpenAlex's and Semantic
Scholar's keyword searches for the studies reviews include, and of the
databases computing reviews need; such a study would say whether the
six databases of \cref{ch:procedure} are enough or which one to add.

\section{The records judged bearing}
\label{sec:databases-hits}

The first table lists the bearing records decided in the question's
session; the second those decided in its companion session of the first
100 places, which holds seven bearing records the first table lacks (the
cited works appear in both).

% Author's note: tables by generate_audit_table on copies of screenlit-q3
% and screenlit-q3-top, bearing rows only, cited = cited in
% doc/bibliography.bib.  Five of the seven bearing records of the -top
% session are pending in screenlit-q3; two (arXiv) are not in it.
% --- scholar LaTeX fragment --------------------------------------------
% Generated by: scholar sessions export -f table
% Session: screenlit-q3
% \input this file from the body of a document whose
% preamble loads:
%   \usepackage{longtable}
%   \usepackage{booktabs}
%   \usepackage{array}
% ----------------------------------------------------------------------

\begingroup\footnotesize
\begin{longtable}{@{}%
  >{\raggedright\arraybackslash}p{0.44\textwidth}c%
  >{\raggedright\arraybackslash}p{0.13\textwidth}%
  >{\raggedright\arraybackslash}p{0.26\textwidth}@{}}
\caption{Records bearing on the claim, Question 3, databases and ranked lists (15 cited, 17 supporting and 23 qualifying or contradicting it, out of 60,049 unique hits; the 158 adjacent hits, the 814 false hits and 59,022 unclassified are not listed). The last column gives the reason for each inclusion; for rows decided by an agent the agent's confidence is given. A row is cited when the text cites the work, whoever decided on its records. Databases searched: OpenAlex (OA; not open access), IEEE Xplore (IEEE), Web of Science (WoS), Scopus (incomplete for 1 of 10 queries), Semantic Scholar (S2; incomplete for 4 of 10 queries), arXiv. Searched 2026-11-23.}\label{tab:screenlit-q3}\\
\toprule
Title & Year & Database & Reason \\
\midrule
\endfirsthead
\multicolumn{4}{@{}l}{\emph{\tablename~\thetable{} (cont.)}}\\
\toprule
Title & Year & Database & Reason \\
\midrule
\endhead
\midrule \multicolumn{4}{r@{}}{\emph{continued on next page}}\\
\endfoot
\bottomrule
\endlastfoot
A Comparison of Citation Sources for Reference and Citation-Based Search in Systematic Literature Reviews & 2022 & WoS, Scopus & \emph{cited}: qualifies or contradicts the claim (0.60) \\
Assessing the coverage of PubMed, Embase, OpenAlex, and Semantic Scholar for automated single-database searches in living guideline evidence surveillance: a case study of the international polycystic ovary syndrome guidelines 2023 & 2025 & OA, WoS, Scopus, S2 & \emph{cited}: supports the claim (0.75) \\
Can electronic search engines optimize screening of search results in systematic reviews: an empirical study & 2006 & OA, Scopus & \emph{cited}: qualifies or contradicts the claim (0.60) \\
Comparing the coverage, recall, and precision of searches for 120 systematic reviews in Embase, MEDLINE, and Google Scholar: a prospective study & 2016 & OA, WoS, Scopus, S2 & \emph{cited}: qualifies or contradicts the claim (0.70) \\
Coverage and similarity of bibliographic databases to find most relevant literature for systematic reviews in education & 2024 & WoS, OA, Scopus, S2 & \emph{cited}: qualifies or contradicts the claim (0.75) \\
Evaluating the proportion of retrieved citations and relevant studies when database selection is restricted to MEDLINE, CENTRAL, and Embase: part 4 of a meta-research study of rapid review methodological trade-offs & 2026 & Scopus, OA & \emph{cited}, pending, unclassified \\
Feasibility and desirability of screening search results from Google Search exhaustively for systematic reviews: A cross-case analysis & 2023 & Scopus & \emph{cited}, pending, unclassified \\
Is the coverage of google scholar enough to be used alone for systematic reviews & 2013 & WoS, Scopus, S2, OA & \emph{cited}: supports the claim (0.60) \\
Optimal database combinations for literature searches in systematic reviews: a prospective exploratory study & 2017 & OA, WoS, Scopus, S2 & \emph{cited}: supports the claim (0.80) \\
Optimizing literature search in systematic reviews - are MEDLINE, EMBASE and CENTRAL enough for identifying effect studies within the area of musculoskeletal disorders? & 2016 & Scopus, OA, WoS, S2 & \emph{cited}: supports the claim (0.60) \\
The contribution of databases to the results of systematic reviews: a cross-sectional study & 2016 & OA, WoS, Scopus & \emph{cited}, pending, unclassified \\
The Impact Of Searching Fewer Databases In Health Technology Assessment Rapid Reviews & 2017 & OA & \emph{cited}, pending, unclassified \\
The Role of Google Scholar in Evidence Reviews and Its Applicability to Grey Literature Searching & 2015 & OA, WoS, Scopus, S2 & \emph{cited}: qualifies or contradicts the claim (0.65) \\
Where you search determines what you find: the effects of bibliographic databases on systematic reviews & 2022 & WoS, OA & \emph{cited}: qualifies or contradicts the claim (0.75) \\
Which academic search systems are suitable for systematic reviews or meta-analyses? Evaluating retrieval qualities of Google Scholar, PubMed, and 26 other resources & 2020 & WoS & \emph{cited}: qualifies or contradicts the claim (0.80) \\
A Comparison of OpenAlex With Scopus and Web of Science for Tracking Scholarly Nursing Literature & 2025 & Scopus, OA, WoS, S2 & supports the claim (0.55) \\
A search of only four key databases would identify most randomized controlled trials of acupuncture: A metaepidemiological study & 2022 & OA, WoS, Scopus, S2 & supports the claim (0.70) \\
Analyzing the Utility of Openalex to Identify Studies for Systematic Reviews: Methods and a Case Study & 2025 & OA, S2 & supports the claim (0.55) \\
Can Google Scholar Serve as a Sufficient Source for Systematic Reviews? Reconsidering the Need for Multiple Database Searches & 2026 & OA & supports the claim (0.40) \\
Database coverage and their use in systematic reviews regarding spinal manipulative therapy: an exploratory study & 2022 & WoS, Scopus, OA, S2 & supports the claim (0.55) \\
Google Scholar as a Sole Literature Source for Meta-analyses in Otolaryngology Head and Neck Surgery & 2023 & Scopus, OA, S2, WoS & supports the claim (0.60) \\
Google Scholar Could Be Used as a Stand-Alone Resource for Systematic Reviews & 2015 & WoS, OA & supports the claim (0.50) \\
Google Scholar Out-Performs Many Subscription Databases when Keyword Searching & 2010 & OA, WoS & supports the claim (0.40) \\
Google Scholar, Microsoft Academic, Scopus, Dimensions, Web of Science, and OpenCitations' COCI: a multidisciplinary comparison of coverage via citations & 2021 & WoS, Scopus & supports the claim (0.45) \\
Guidelines for the Search Strategy to Update Systematic Literature Reviews in Software Engineering & 2020 & arXiv, WoS, Scopus, S2 & supports the claim (0.50) \\
Knowing When to Stop: Final Results vs. Work Involved in Systematic Review Database Searching & 2016 & OA & supports the claim (0.55) \\
On the performance of hybrid search strategies for systematic literature reviews in software engineering & 2020 & OA, arXiv, WoS, Scopus, S2 & supports the claim (0.55) \\
Reassessing the Role of Google Scholar in PRISMA-Informed Systematic Reviews Through a Critical Analysis of Influential Research Identifying Its Limitations and Empirical Evidence & 2026 & OA, WoS, Scopus, S2 & supports the claim (0.50) \\
Redefining searching in non-medical sciences systematic reviews: The ascendance of Google Scholar as the primary database & 2025 & Scopus, OA, WoS, S2 & supports the claim (0.50) \\
Reference coverage analysis of OpenAlex compared to Web of Science and Scopus & 2025 & OA, WoS, Scopus, S2 & supports the claim (0.50) \\
Sources Used To Find Studies For Systematic Reviews Of Economic Evaluations & 2017 & OA, S2 & supports the claim (0.50) \\
The Utility of Google Scholar When Searching Geographical Literature: Comparison With Three Commercial Bibliographic Databases & 2015 & OA, WoS, Scopus, S2 & supports the claim (0.40) \\
A comparison of the performance of seven key bibliographic databases in identifying all relevant systematic reviews of interventions for hypertension & 2016 & OA, WoS, Scopus, S2 & qualifies or contradicts the claim (0.55) \\
A systematic review case study of urgent and emergency care configuration found citation searching of Web of Science and Google Scholar of similar value & 2022 & OA, WoS, Scopus, S2 & qualifies or contradicts the claim (0.55) \\
A Systematic Review on the Sufficiency of PubMed and Google Scholar for Biosciences & 2020 & OA, S2 & qualifies or contradicts the claim (0.55) \\
An alternative screening approach for Google Search identifies an accurate and manageable number of results for a systematic review (case study) & 2021 & OA, WoS, Scopus, S2 & qualifies or contradicts the claim (0.55) \\
An analysis of the suitability of OpenAlex for bibliometric analyses & 2024 & OA, S2 & qualifies or contradicts the claim (0.50) \\
Beyond Google Scholar, Scopus, and Web of Science: An evaluation of the backward and forward citation coverage of 59 databases' citation indices & 2024 & Scopus, OA, WoS, S2 & qualifies or contradicts the claim (0.50) \\
Can we prioritise which databases to search? A case study using a systematic review of frozen shoulder management & 2012 & OA, WoS, Scopus, S2 & qualifies or contradicts the claim (0.55) \\
Citation searching: a systematic review case study of multiple risk behaviour interventions & 2014 & OA, WoS, Scopus, S2 & qualifies or contradicts the claim (0.50) \\
Coverage of the bibliographic databases in mental health research & 2009 & OA, WoS, Scopus, S2 & qualifies or contradicts the claim (0.45) \\
Evaluating the linguistic coverage of OpenAlex : An assessment of metadata accuracy and completeness & 2025 & OA, WoS, Scopus, S2 & qualifies or contradicts the claim (0.45) \\
Google Scholar as replacement for systematic literature searches: good relative recall and precision are not enough & 2013 & OA, WoS, Scopus, S2 & qualifies or contradicts the claim (0.50) \\
Google Scholar is not enough to be used alone for systematic reviews & 2013 & OA & qualifies or contradicts the claim (0.65) \\
How many sources are needed? The effects of bibliographic databases on systematic review outcomes & 2022 & WoS, OA, Scopus, S2, IEEE & qualifies or contradicts the claim (0.70) \\
Identifying Relevant Controlled Clinical Trials for Systematic Reviews Requires Searching Multiple Resources -- and, Even Then, Comprehensiveness is Questionable & 2006 & OA, S2 & qualifies or contradicts the claim (0.55) \\
Large coverage fluctuations in Google Scholar: A case study & 2021 & Scopus, WoS, OA, arXiv, S2 & qualifies or contradicts the claim (0.50) \\
Overlaps of multiple database retrieval and citation tracking in dementia care research: a methodological study & 2021 & OA, WoS, Scopus, S2 & qualifies or contradicts the claim (0.50) \\
Search where you will find most: Comparing the disciplinary coverage of 56 bibliographic databases & 2022 & OA, WoS, Scopus, S2 & qualifies or contradicts the claim (0.50) \\
Searching for qualitative health research required several databases and alternative search strategies: a study of coverage in bibliographic databases & 2019 & OA, WoS, Scopus, S2 & qualifies or contradicts the claim (0.60) \\
Searching the Veterinary Literature: A Comparison of the Coverage of Veterinary Journals by Nine Bibliographic Databases & 2012 & OA, WoS, Scopus, S2 & qualifies or contradicts the claim (0.45) \\
Understanding discrepancies in the coverage of OpenAlex : The case of China & 2025 & OA, WoS, Scopus, S2 & qualifies or contradicts the claim (0.45) \\
Using Conventional Bibliographic Databases for Social Science Research: Web of Science and Scopus are not the Only Options & 2021 & OA, Scopus, S2 & qualifies or contradicts the claim (0.40) \\
Variation in number of hits for complex searches in Google Scholar & 2016 & WoS, OA & qualifies or contradicts the claim (0.45) \\
Web Search Engines - Not Yet a Reliable Replacement for Bibliographic Databases & 2018 & OA, WoS, Scopus, S2 & qualifies or contradicts the claim (0.60) \\
\end{longtable}
\endgroup

% --- scholar LaTeX fragment --------------------------------------------
% Generated by: scholar sessions export -f table
% Session: screenlit-q3-top
% \input this file from the body of a document whose
% preamble loads:
%   \usepackage{longtable}
%   \usepackage{booktabs}
%   \usepackage{array}
% ----------------------------------------------------------------------

\begingroup\footnotesize
\begin{longtable}{@{}%
  >{\raggedright\arraybackslash}p{0.44\textwidth}c%
  >{\raggedright\arraybackslash}p{0.13\textwidth}%
  >{\raggedright\arraybackslash}p{0.26\textwidth}@{}}
\caption{Records bearing on the claim, Question 3, records decided in the companion session of the first 100 places (9 cited, 1 supporting and 6 qualifying or contradicting it, out of 979 unique hits; the 11 adjacent hits, the 85 false hits and 867 unclassified are not listed). The last column gives the reason for each inclusion; for rows decided by an agent the agent's confidence is given. A row is cited when the text cites the work, whoever decided on its records. Databases searched: OpenAlex (OA; not open access; capped at the limit for 6 of 10 queries), arXiv. Searched 2026-11-23.}\label{tab:screenlit-q3-top}\\
\toprule
Title & Year & Database & Reason \\
\midrule
\endfirsthead
\multicolumn{4}{@{}l}{\emph{\tablename~\thetable{} (cont.)}}\\
\toprule
Title & Year & Database & Reason \\
\midrule
\endhead
\midrule \multicolumn{4}{r@{}}{\emph{continued on next page}}\\
\endfoot
\bottomrule
\endlastfoot
Assessing the coverage of PubMed, Embase, OpenAlex, and Semantic Scholar for automated single-database searches in living guideline evidence surveillance: a case study of the international polycystic ovary syndrome guidelines 2023 & 2025 & OA & \emph{cited}, pending, unclassified \\
Can electronic search engines optimize screening of search results in systematic reviews: an empirical study & 2006 & OA & \emph{cited}, pending, unclassified \\
Comparing the coverage, recall, and precision of searches for 120 systematic reviews in Embase, MEDLINE, and Google Scholar: a prospective study & 2016 & OA & \emph{cited}, pending, unclassified \\
Coverage and similarity of bibliographic databases to find most relevant literature for systematic reviews in education & 2023 & OA & \emph{cited}, pending, unclassified \\
Is the coverage of google scholar enough to be used alone for systematic reviews & 2013 & OA & \emph{cited}, pending, unclassified \\
Optimal database combinations for literature searches in systematic reviews: a prospective exploratory study & 2017 & OA & \emph{cited}, pending, unclassified \\
Optimizing literature search in systematic reviews -- are MEDLINE, EMBASE and CENTRAL enough for identifying effect studies within the area of musculoskeletal disorders? & 2016 & OA & \emph{cited}, pending, unclassified \\
The Role of Google Scholar in Evidence Reviews and Its Applicability to Grey Literature Searching & 2015 & OA & \emph{cited}, pending, unclassified \\
Where you search determines what you find: the effects of bibliographic databases on systematic reviews & 2021 & OA & \emph{cited}, pending, unclassified \\
Complementary approaches to searching MEDLINE may be sufficient for updating systematic reviews & 2016 & OA & supports the claim (0.50) \\
Cochrane Systematic Reviews in Acupuncture: Methodological Diversity in Database Searching & 2005 & OA & qualifies or contradicts the claim (0.45) \\
Combination of multiple databases is necessary for a valid systematic review & 2014 & OA & qualifies or contradicts the claim (0.60) \\
Evaluating the linguistic coverage of OpenAlex : An assessment of metadata accuracy and completeness & 2025 & OA, arXiv & qualifies or contradicts the claim (0.40) \\
Literature search in systematic reviews: How much is good enough? & 2023 & OA & qualifies or contradicts the claim (0.60) \\
Searching Additional Databases Except PubMed Are Necessary for a Systematic Review & 2008 & OA & qualifies or contradicts the claim (0.55) \\
Understanding discrepancies in the coverage of OpenAlex : The case of China & 2025 & OA, arXiv & qualifies or contradicts the claim (0.40) \\
\end{longtable}
\endgroup

